%Data institusi

%\input{Dozentdaten}

%\renewcommand{\fachbereich}{Fachbereich}

%\renewcommand{\dozent}{Prof. Dr. . }

%Data ujian

\renewcommand{\klausurgebiet}{ }

\renewcommand{\klausurtyp}{ }

\renewcommand{\klausurdatum}{ . 20}

\klausurvorspann {\fachbereich} {\klausurdatum} {\dozent} {\klausurgebiet} {\klausurtyp}

%Data untuk tabel nilai berikut

\renewcommand{\aeins}{  3  }

\renewcommand{\azwei}{  3   }

\renewcommand{\adrei}{  0  }

\renewcommand{\avier}{  0  }

\renewcommand{\afuenf}{  6  }

\renewcommand{\asechs}{  6  }

\renewcommand{\asieben}{  0  }

\renewcommand{\aacht}{  6  }

\renewcommand{\aneun}{  11  }

\renewcommand{\azehn}{  5  }

\renewcommand{\aelf}{  5  }

\renewcommand{\azwoelf}{  3  }

\renewcommand{\adreizehn}{  7   }

\renewcommand{\avierzehn}{  2  }

\renewcommand{\afuenfzehn}{  0  }

\renewcommand{\asechzehn}{  6  }

\renewcommand{\asiebzehn}{  63  }

\renewcommand{\aachtzehn}{    }

\renewcommand{\aneunzehn}{    }

\renewcommand{\azwanzig}{    }

\renewcommand{\aeinundzwanzig}{    }

\renewcommand{\azweiundzwanzig}{    }

\renewcommand{\adreiundzwanzig}{    }

\renewcommand{\avierundzwanzig}{    }

\renewcommand{\afuenfundzwanzig}{    }

\renewcommand{\asechsundzwanzig}{    }

\punktetabellesechzehn

\klausurnote

\newpage

\setcounter{section}{K}

\inputaufgabegibtloesung
{3}
{

Definisikan istilah-istilah berikut
\zusatzklammer {yang dicetak miring} {} {.}
\aufzaehlungsechs{Suatu
\stichwort {kurva geodesik} {}
pada hiperpermukaan diferensiabel

\mavergleichskette
{\vergleichskette
{ Y
}
{ \subseteq }{ \R^n
}
{  }{
}
{    }{
}
{   }{
}
}
{}{}{.}

}{
\stichwort {Himpunan rotasi} {}
\zusatzklammer {mengelilingi sumbu $x$} {} {}
dari suatu subhimpunan

\mavergleichskette
{\vergleichskette
{ T
}
{ \subseteq }{  \R \times \R_{\geq 0}
}
{  }{
}
{    }{
}
{   }{
}
}
{}{}{.}

}{
\stichwort {Pemetaan tangen} {}
\maabbdisp {T (\varphi)} {TM} {TN
} {}
dari
\definitionsverweis {pemetaan diferensiabel}{}{}
\maabbdisp {\varphi} { M } { N
} {}
antara
\definitionsverweis {manifold-manifold diferensiabel}{}{}
\mathkor {} {M} {dan} {N} {.}

}{Suatu perubahan peta yang
\stichwort {mempertahankan orientasi} {}
pada
\definitionsverweis {manifold diferensiabel}{}{}
$M$.

}{
\stichwort {Bentuk diferensial tarikan balik} {} $\varphi^* \omega$ dari suatu bentuk diferensial
\mathl{\omega \in
{ \mathcal E }^{ k } (  M   )}{} terhadap suatu pemetaan terdiferensialkan secara kontinu
\maabb {\varphi} { L } { M
} {}
antara dua manifold diferensiabel
\mathkor {} {L} {dan} {M} {.}

}{Suatu koneksi linear yang
\stichwort {metrik} {}
pada bundel tangen suatu manifold Riemann $M$.
}

}
{} {}

\inputaufgabegibtloesung
{3}
{

Nyatakan teorema-teorema berikut.
\aufzaehlungdrei{
\stichwort {Teorema tentang solusi persamaan diferensial biasa pada hiperpermukaan diferensiabel} {}

\mavergleichskette
{\vergleichskette
{  Y
}
{ \subseteq }{  \R^n
}
{  }{
}
{    }{
}
{   }{
}
}
{}{}{.}}{Rumus untuk menghitung volume kanonik suatu manifold Riemann dalam sebuah peta.}{Teorema tentang partisi kesatuan.}

}
{} {}

\inputaufgabe
{0}
{

}
{} {}

\inputaufgabe
{0}
{

}
{} {}

\inputaufgabegibtloesung
{6}
{

Untuk permukaan $Y$ yang diberikan oleh

\mavergleichskettedisp
{\vergleichskette
{ x^2+y^2+2z^2
}
{ =} { 4
}
{ } {
}
{ } {
}
{ } {
}
}
{}{}{}
dan titik

\mavergleichskettedisp
{\vergleichskette
{ P
}
{ =} { \left(  1  , \,   1  , \,   1                 \right)
}
{ \in} { Y
}
{ } {
}
{ } {
}
}
{}{}{,}
tentukan suatu matriks diagonal bagi
\definitionsverweis {pemetaan Weingarten}{}{}
$L_P$.

}
{} {}

\inputaufgabegibtloesung
{6}
{

Buktikan teorema isometri untuk transpor paralel pada hiperpermukaan

\mavergleichskette
{\vergleichskette
{ Y
}
{ \subseteq }{ \R^n
}
{  }{
}
{    }{
}
{   }{
}
}
{}{}{.}

}
{} {}

\inputaufgabe
{0}
{

}
{} {}

\inputaufgabegibtloesung
{6}
{

Misalkan $M$ suatu manifold topologis kompak berdimensi $d$,

\mavergleichskette
{\vergleichskette
{ d
}
{ \geq }{ 1
}
{  }{
}
{    }{
}
{   }{
}
}
{}{}{.}
Tunjukkan bahwa terdapat suatu subhimpunan terbuka terbatas

\mavergleichskette
{\vergleichskette
{ U
}
{ \subseteq }{ \R^d
}
{  }{
}
{    }{
}
{   }{
}
}
{}{}{}
dan suatu pemetaan kontinu surjektif
\maabbdisp {\varphi} { U } { M
} {}
.

}
{} {}

\inputaufgabegibtloesung
{11 (1+3+1+2+4)}
{

Kita tinjau himpunan

\mavergleichskettedisp
{\vergleichskette
{N
}
{ =} { \left\{ A= \begin{pmatrix}  a   &   b   \\  c  & d \end{pmatrix}   \mid   A \text{ nilpoten}         \right\}
}
{ \subseteq} {\R^4
}
{ } {
}
{ } {
}
}
{}{}{}
dari
$2\times 2$-\definitionsverweis {matriks}{}{}
real nilpoten, serta himpunan

\mavergleichskettedisp
{\vergleichskette
{M
}
{ =} { N \setminus \{ \begin{pmatrix}  0   &   0  \\  0  &  0 \end{pmatrix}  \}
}
{ } {
}
{ } {
}
{ } {
}
}
{}{}{.}

a) Apakah $N$ terhubung?

b) Tunjukkan bahwa $M$ suatu
\definitionsverweis {submanifold tertutup}{}{}
dari suatu subhimpunan terbuka

\mavergleichskette
{\vergleichskette
{ G
}
{ \subseteq }{ \R^4
}
{  }{
}
{    }{
}
{   }{
}
}
{}{}{.}

c) Tentukan
\definitionsverweis {dimensi}{}{}
$M$.

d) Apakah $M$ terhubung?

e) Tutupi
\mathl{M}{} dengan peta-peta topologis eksplisit.

}
{} {}

\inputaufgabe
{5}
{

Katakan sesuatu yang cerdas tentang pita Möbius.

}
{} {}

\inputaufgabegibtloesung
{5}
{

Kita tinjau pemetaan
\maabbeledisp {\varphi} { \R^3 } { \R^2
} { (x,y,z) } { \left(  xyz^2,xy-z^3  \right)
} {.}
Hitung matriks pemetaan
\maabbdisp {\bigwedge^2 T_P (\varphi)} { \bigwedge^2 T_P \R^3 } { \bigwedge^2 T_{\varphi(P)}  \R^2
} {}
di titik

\mavergleichskette
{\vergleichskette
{ P
}
{ = }{ (1,3,5)
}
{  }{
}
{    }{
}
{   }{
}
}
{}{}{}
terhadap suatu basis yang sesuai.

}
{} {}

\inputaufgabegibtloesung
{3}
{

Misalkan

\mavergleichskette
{\vergleichskette
{ W
}
{ \subseteq }{ \R^m
}
{  }{
}
{    }{
}
{   }{
}
}
{}{}{}
dan

\mavergleichskette
{\vergleichskette
{ U
}
{ \subseteq }{ \R^n
}
{  }{
}
{    }{
}
{   }{
}
}
{}{}{}
merupakan
\definitionsverweis {subhimpunan-subhimpunan terbuka}{}{}
dan misalkan
\maabbdisp {\psi} { W } { U
} {}
suatu
\definitionsverweis {pemetaan}{}{}
yang
\definitionsverweis {terdiferensialkan secara kontinu}{}{.}
Misalkan
\maabbdisp {f} { U } {\R
} {}
suatu fungsi terdiferensialkan secara kontinu. Simpulkan dari
aturan rantai
bahwa

\mavergleichskettedisp
{\vergleichskette
{ d (\psi^*f)
}
{ =} { \psi^*(df)
}
{ } {
}
{ } {
}
{ } {
}
}
{}{}{,}
dengan $\psi^*$ menyatakan
\definitionsverweis {penarikan balik bentuk diferensial}{}{.}

}
{} {}

\inputaufgabegibtloesung
{7}
{

Misalkan $M$ suatu
\definitionsverweis {manifold diferensiabel}{}{}
dengan
\definitionsverweis {basis terhitung bagi topologinya}{}{}
dan misalkan $\omega$ suatu
\definitionsverweis {bentuk volume positif}{}{}
pada $M$. Misalkan
\maabbdisp {\varphi} { L } { M
} {}
suatu
\definitionsverweis {difeomorfisme}{}{}
dengan manifold $L$ dan

\mavergleichskette
{\vergleichskette
{ S
}
{ \subseteq }{ L
}
{  }{
}
{    }{
}
{   }{
}
}
{}{}{}
suatu
\definitionsverweis {subhimpunan terukur}{}{.}
Tunjukkan bahwa

\mavergleichskettedisp
{\vergleichskette
{ \int_S \varphi^* \omega
}
{ =} { \int_{ \varphi(S)} \omega
}
{ } {
}
{ } {
}
{ } {
}
}
{}{}{.}

}
{} {}

\inputaufgabegibtloesung
{2}
{

Pada
\definitionsverweis {bundel vektor trivial}{}{}
\maabbdisp {p} {\R^2 \times \R} { \R^2
} {}
ber-
\definitionsverweis {rank}{}{}
$1$ di atas $\R^2$, kita tinjau
\definitionsverweis {koneksi linear}{}{,}
yang diberikan oleh
\definitionsverweis {simbol-simbol Christoffel}{}{}

\mavergleichskette
{\vergleichskette
{ \Gamma_1 (x,y)
}
{ = }{ x^5-x^2y^2+3y^4
}
{  }{
}
{    }{
}
{   }{
}
}
{}{}{}
dan

\mavergleichskette
{\vergleichskette
{ \Gamma_2 (x,y)
}
{ = }{ 7x^3 +xy^2-4y^5
}
{  }{
}
{    }{
}
{   }{
}
}
{}{}{.}
Hitung
\definitionsverweis {operator kelengkungan}{}{}
$R \left(  \partial_1, \partial_2  \right)$.

}
{} {}

\inputaufgabe
{0}
{

}
{} {}

\inputaufgabegibtloesung
{6}
{

Buktikan teorema tentang retraksi ke batas pada manifold.

}
{} {}
